% Vietnamese translation for OI Public Library
% Source-PDF: IOI中国国家候选队论文/2002/张一飞.pdf
\documentclass[11pt]{article}
\usepackage{oipl-vn}

\title{Từ nhận thức cảm tính đến nhận thức lý tính: phân tích quá trình giải một lớp trò chơi tổ hợp}
\author{Zhang Yifei}
\date{}

\newcommand{\nimadd}{\oplus}
\newcommand{\val}[1]{\mathop{\#}\!#1}
\newcommand{\moves}[1]{\text{\$}(#1)}
\newcommand{\Nat}{\mathbb{N}}

\begin{document}
\maketitle

\origtitle{From perceptual understanding to rational understanding: analyzing the solution process of a class of combinatorial games}
\sourcepath{IOI China candidate team papers / 2002 / Zhang Yifei.pdf}

\tableofcontents

\section{Trò chơi}

\paragraph{Trò chơi A.}
Hai người, Giáp và Ất, đối mặt với một số đống sỏi. Số sỏi trong mỗi đống có
thể được chọn tùy ý. Chẳng hạn, một vị trí ban đầu có \(n=3\) đống:
\[
  a_1=3,\qquad a_2=3,\qquad a_3=1.
\]
Hai người lần lượt lấy sỏi theo các quy tắc sau:
\begin{itemize}
  \item mỗi nước đi phải lấy ít nhất một viên sỏi;
  \item mỗi nước đi chỉ được lấy một phần hoặc toàn bộ sỏi từ một đống;
  \item ai không thể đi đúng quy tắc thì thua.
\end{itemize}
Ta có thể biểu diễn ví dụ trên bằng sơ đồ văn bản:
\[
\begin{array}{ccc}
\text{Đống thứ nhất: }a_1=3&
\text{Đống thứ hai: }a_2=3&
\text{Đống thứ ba: }a_3=1
\end{array}
\]

\paragraph{Trò chơi B.}
Giáp và Ất thỏa thuận trước một số \(m\), và mỗi lần không được lấy quá \(m\)
viên sỏi. Các quy tắc còn lại giống trò chơi A.

Điều ta quan tâm là: với một vị trí ban đầu, người đi trước, tức Giáp, có
chiến lược thắng chắc hay người đi sau, tức Ất, có chiến lược thắng chắc.
Sau đây ta bắt đầu từ các trường hợp đơn giản để nghiên cứu một số tính chất
của trò chơi này.

\section{Bắt đầu từ trường hợp đơn giản}

Dùng một bộ \(n\) phần tử \((a_1,a_2,\ldots,a_n)\) để mô tả một vị trí trong
quá trình chơi. Ví dụ, vị trí ở trên được mô tả bởi bộ
\((3,3,1)\). Đổi thứ tự các số trong bộ này vẫn biểu diễn cùng một vị trí;
chẳng hạn \((3,3,1)\) và \((1,3,3)\) được xem là như nhau.

Nếu vị trí ban đầu chỉ có một đống sỏi thì Giáp có chiến lược thắng chắc:
Giáp lấy hết đống đó trong một lần, khiến Ất không còn sỏi để lấy.

Nếu vị trí ban đầu có hai đống sỏi và hai đống có số sỏi bằng nhau thì Ất có
chiến lược thắng chắc. Vì có hai đống nên Giáp không thể lấy hết ngay trong
một lần. Nếu Giáp lấy một số sỏi trong một đống, Ất lấy đúng từng ấy sỏi ở
đống kia. Theo tính đối xứng, sau khi Giáp đi, Ất luôn còn nước đi. Tổng số
sỏi giảm dần, nên cuối cùng chắc chắn Giáp là người không còn nước đi. Vì
thế vị trí \((1)\) là vị trí thắng cho người đi trước, còn \((3,3)\) là vị
trí thua cho người đi trước.

\paragraph{Phép cộng vị trí.}
Định nghĩa
\[
  (a_1,a_2,\ldots,a_n)+(b_1,b_2,\ldots,b_m)
  =(a_1,a_2,\ldots,a_n,b_1,b_2,\ldots,b_m).
\]
Ví dụ,
\[
  (3)+(3)+(1)=(3,3)+(1)=(3,3,1).
\]
Với các vị trí \(A,B,S\), nếu \(S=A+B\), ta nói \(S\) có thể phân tích thành
hai vị trí con \(A\) và \(B\). Chẳng hạn, \((3,3,1)\) có thể phân tích thành
\((3,3)\) và \((1)\).

Nếu vị trí ban đầu có thể phân thành hai vị trí con giống nhau thì Ất có
chiến lược thắng chắc. Thật vậy, giả sử \(S=A+A\). Hãy tưởng tượng có hai
bàn, trên mỗi bàn đặt một vị trí \(A\). Nếu Giáp lấy sỏi ở một bàn thì Ất
lấy đối xứng ở bàn kia. Nhờ đối xứng, sau khi Giáp đi, Ất luôn còn nước đi;
tổng số sỏi giảm dần, nên cuối cùng Giáp không thể đi. Ví dụ,
\[
  (2,2,5,5,5,5,7,7)=(2,5,5,7)+(2,5,5,7)
\]
là vị trí thua cho người đi trước.

Nếu người đi trước có chiến lược thắng chắc ở vị trí \(S\), ta nói \(S\) là
\term{vị trí thắng}. Nếu người đi sau có chiến lược thắng chắc, ta nói \(S\)
là \term{vị trí thua}. Ví dụ, với
\[
  A=(1),\qquad B=(3,3),\qquad C=(2,2,5,5,5,5,7,7),
\]
ta có \(A\) thắng, còn \(B\) và \(C\) thua. Vấn đề của ta là tìm cách phán
đoán thắng thua của một vị trí.

Hai tính chất cơ bản là:
\begin{itemize}
  \item nếu \(S\) thắng thì tồn tại một cách lấy sỏi \(S\to T\) sao cho
  \(T\) thua;
  \item nếu \(S\) thua thì với mọi cách lấy sỏi \(S\to T\), vị trí \(T\)
  thắng.
\end{itemize}

\paragraph{Lý thuyết phân tích vị trí.}
Giả sử \(S=A+B\). Khi đó:
\begin{itemize}
  \item nếu \(A\) và \(B\) một thắng một thua thì \(S\) thắng. Không mất tính
  tổng quát, giả sử \(A\) thắng và \(B\) thua. Hãy tưởng tượng hai bàn mang
  hai vị trí \(A\) và \(B\). Vì \(A\) thắng, Giáp có thể bảo đảm lấy viên sỏi
  cuối cùng trên bàn \(A\). Đồng thời, Giáp cũng có thể bảo đảm người đi
  nước đầu tiên trên bàn \(B\) là Ất. Vì \(B\) thua, Giáp cũng bảo đảm lấy
  viên sỏi cuối cùng trên bàn \(B\). Vậy Giáp lấy được viên cuối cùng trên
  cả hai bàn.

  \item nếu \(A\) thua và \(B\) thua thì \(S\) thua. Dù Giáp đi trước trong
  \(A\) hay trong \(B\), vị trí còn lại cho Ất luôn là tổng của một vị trí
  thắng và một vị trí thua, tức là vị trí thắng.

  \item nếu \(B\) thua thì thắng thua của \(S\) giống thắng thua của \(A\);

  \item nếu \(A\) thắng và \(B\) thắng thì \(S\) có khi thắng, có khi thua.
\end{itemize}

Suy ra nếu \(S=A+C+C\) thì thắng thua của \(S\) giống thắng thua của \(A\).
Lấy \(B=C+C\), ta có \(S=A+B\) và \(B\) thua. Ví dụ,
\[
  (3,3,1)=(3)+(3)+(1)
\]
có thắng thua giống \((1)\), nên \((3,3,1)\) là vị trí thắng.

Gọi vị trí không có viên sỏi nào là \term{vị trí rỗng}; vị trí rỗng là vị trí
thua.

Nếu trong vị trí \(S\) có hai đống sỏi bằng nhau, và \(T\) là vị trí thu
được sau khi bỏ hai đống ấy khỏi \(S\), ta nói \(S\) có thể \term{đơn giản
hóa} thành \(T\). Ví dụ,
\[
  (2,2,2,7,9,9)\to (2,2,2,7)\to (2,7).
\]
Một vị trí có cùng thắng thua với vị trí sau khi đơn giản hóa. Vị trí không
thể đơn giản hóa nữa gọi là \term{vị trí tối giản}. Trong vị trí tối giản
không có hai đống bằng nhau, nên có thể dùng một tập hợp để biểu diễn nó.
Chẳng hạn, vị trí tối giản \((2,7)\) có thể biểu diễn bằng tập \(\{2,7\}\).

Nếu chỉ quan tâm thắng thua, một vị trí có thể được mô tả bằng một tập hợp.
Ví dụ, \((3,3,1)\) có thể được mô tả bằng \(\{1\}\). Nếu giải trò chơi bằng
tìm kiếm trên cây trò chơi, dùng tập hợp thay cho bộ nhiều phần tử sẽ giảm
số vị trí phải xét, nhưng độ phức tạp thời gian vẫn rất cao. Liệu còn có thể
đơn giản hóa cách biểu diễn vị trí hơn nữa không?

\section{Tương tự và liên tưởng}

Trong bài này, \term{phép cộng nhị phân} là phép cộng nhị phân không nhớ, tức
phép xor logic. Với từng bit:
\[
  1\nimadd 0=1,\qquad 0\nimadd 1=1,\qquad
  0\nimadd 0=0,\qquad 1\nimadd 1=0.
\]

Ta so sánh phép cộng nhị phân với phép cộng vị trí. Dùng chữ hoa \(A,B\) cho
vị trí, chữ thường \(a,b\) cho số nhị phân:
\begin{itemize}
  \item nếu \(A\) và \(B\) giống nhau thì \(A+B\) thua; nếu \(a=b\) thì
  \(a\nimadd b=0\);
  \item nếu \(A\) thắng, \(B\) thua thì \(A+B\) thắng; nếu \(a\ne0\) và
  \(b=0\) thì \(a\nimadd b\ne0\);
  \item nếu \(B\) thắng, \(A\) thua thì \(A+B\) thắng; nếu \(b\ne0\) và
  \(a=0\) thì \(a\nimadd b\ne0\);
  \item nếu \(A\) thua, \(B\) thua thì \(A+B\) thua; nếu \(a=0\) và \(b=0\)
  thì \(a\nimadd b=0\);
  \item nếu \(A\) thắng, \(B\) thắng thì \(A+B\) có khi thắng, có khi thua;
  nếu \(a\ne0\) và \(b\ne0\) thì \(a\nimadd b\) có khi khác \(0\), có khi
  bằng \(0\).
\end{itemize}

Nếu dùng số nhị phân \(s\) để biểu diễn thắng thua của vị trí \(S\), trong đó
\(S\) thắng ứng với \(s\ne0\) và \(S\) thua ứng với \(s=0\), thì phép cộng vị
trí có đúng các tính chất của phép cộng nhị phân không nhớ. Ta có thể dùng
một số nhị phân để biểu diễn một vị trí không?

Dùng ký hiệu \(\val{S}\) để chỉ số nhị phân ứng với vị trí \(S\). Nếu \(S\)
chỉ có một đống sỏi, ta dùng số nhị phân ứng với số sỏi trong đống đó để biểu
diễn \(S\). Ví dụ,
\[
  \val{(5)}=5=(101)_2.
\]
Nếu \(S=A+B\) thì
\[
  \val{S}=\val{A}\nimadd\val{B}.
\]
Vì vậy,
\[
  \val{(3,3)}=\val{(3)}\nimadd\val{(3)}=(11)_2\nimadd(11)_2=0,
\]
và
\[
  \val{(3,3,1)}=\val{(3,3)}\nimadd\val{(1)}=0\nimadd1=1.
\]

Định nghĩa hàm \(f\): nếu vị trí \(S\) chỉ có một đống sỏi, \(S=\{a_1\}\),
thì \(f(a_1)=\val{S}\). Với trò chơi A, \(\val{(5)}=(101)_2\), nên
\(f(5)=(101)_2\); nói cách khác \(f(x)=x\).

Nếu
\[
  S=(a_1,a_2,\ldots,a_n)=(a_1)+(a_2)+\cdots+(a_n),
\]
thì
\[
  \val{S}=f(a_1)\nimadd f(a_2)\nimadd\cdots\nimadd f(a_n).
\]
Ví dụ,
\[
  \val{(3,3,1)}
  =f(3)\nimadd f(3)\nimadd f(1)
  =(11)_2\nimadd(11)_2\nimadd1=1.
\]
Ta kỳ vọng rằng với vị trí \(S\), nếu \(\val{S}=0\) thì \(S\) thua, còn nếu
\(\val{S}\ne0\) thì \(S\) thắng.

\section{Chứng minh}

Trước hết nhắc hai tính chất của phép cộng nhị phân không nhớ:
\begin{itemize}
  \item với hai số nhị phân \(a,b\), \(a\nimadd b=0\) khi và chỉ khi \(a=b\);
  \item với ba số nhị phân \(a,b,s\), nếu \(a\nimadd b=s\) thì
  \(a=b\nimadd s\).
\end{itemize}
Ví dụ:
\[
\begin{array}{cc}
(1011)_2\nimadd(1011)_2=0, &
(0011)_2\nimadd(1010)_2=(1001)_2.
\end{array}
\]

Nếu
\[
  a_1\nimadd a_2\nimadd\cdots\nimadd a_n=p\ne0,
\]
thì tồn tại \(k\) sao cho \(a_k\nimadd p<a_k\). Thật vậy, vì \(p\ne0\), bit
cao nhất của \(p\) bằng \(1\). Gọi đó là bit thứ \(q\). Phải có ít nhất một
\(a_k\) cũng có bit thứ \(q\) bằng \(1\), nếu không bit thứ \(q\) của tổng
xor không thể bằng \(1\). Khi đó bit thứ \(q\) của \(a_k\nimadd p\) bằng
\(0\), còn các bit cao hơn không đổi, nên \(a_k\nimadd p<a_k\).

Nếu \(\val{S}=0\), thì dù người đi trước lấy sỏi như thế nào \(S\to T\), ta
đều có \(\val{T}\ne0\). Giả sử người đi trước chỉ thay đổi đống thứ nhất, lấy
\(x>0\) viên. Với
\[
  S=(a_1,a_2,\ldots,a_n),\qquad
  T=(a_1-x,a_2,\ldots,a_n),
\]
ta có
\[
  \val{S}=f(a_1)\nimadd\val{(a_2,\ldots,a_n)}=0,
\]
nên
\[
  f(a_1)=\val{(a_2,\ldots,a_n)}.
\]
Do đó
\[
\begin{aligned}
  \val{T}
  &=f(a_1-x)\nimadd\val{(a_2,\ldots,a_n)}\\
  &=f(a_1-x)\nimadd f(a_1).
\end{aligned}
\]
Vì trong trò chơi A, \(f(y)=y\), từ \(x>0\) suy ra \(f(a_1-x)\ne f(a_1)\),
nên \(\val{T}\ne0\).

Ngược lại, nếu \(\val{S}\ne0\), người đi trước luôn có một nước đi \(S\to T\)
sao cho \(\val{T}=0\). Đặt
\[
  p=\val{S}=f(a_1)\nimadd f(a_2)\nimadd\cdots\nimadd f(a_n).
\]
Vì \(p\ne0\), tồn tại \(k\) sao cho \(f(a_k)\nimadd p<f(a_k)\). Không mất
tính tổng quát, lấy \(k=1\), và đặt
\[
  x=f(a_1)\nimadd p.
\]
Người đi trước giảm số sỏi của đống thứ nhất từ \(a_1\) xuống \(x\). Khi đó
\[
  p=f(a_1)\nimadd\val{(a_2,\ldots,a_n)}
\]
nên
\[
  \val{(a_2,\ldots,a_n)}=f(a_1)\nimadd p=x.
\]
Vì vậy vị trí mới \(T=(x,a_2,\ldots,a_n)\) thỏa
\[
  \val{T}=f(x)\nimadd\val{(a_2,\ldots,a_n)}=x\nimadd x=0.
\]
Vị trí rỗng cũng có giá trị \(0\). Từ hai chiều trên, với trò chơi A:
\[
  \val{S}=0 \Longleftrightarrow S\text{ thua},\qquad
  \val{S}\ne0 \Longleftrightarrow S\text{ thắng}.
\]
Ví dụ,
\[
  \val{(1,2,3)}=01\nimadd10\nimadd11=0,
\]
nên \((1,2,3)\) thua; còn
\[
  \val{(1,2,3,4)}=001\nimadd010\nimadd011\nimadd100=100,
\]
nên \((1,2,3,4)\) thắng.

Với trò chơi A, với mọi vị trí ban đầu
\[
  S=(a_1,a_2,\ldots,a_n),
\]
ta xem các \(a_i\) là số nhị phân và đặt
\[
  \val{S}=a_1\nimadd a_2\nimadd\cdots\nimadd a_n.
\]
Nếu \(\val{S}\ne0\) thì người đi trước, Giáp, có chiến lược thắng chắc; nếu
\(\val{S}=0\) thì người đi sau, Ất, có chiến lược thắng chắc.

Kết luận này mở rộng sang trò chơi B. Đặt
\[
  f(x)=x\bmod(m+1),
\]
và xem giá trị của \(f\) là một số nhị phân. Với mọi vị trí
\[
  S=(a_1,a_2,\ldots,a_n),
\]
đặt
\[
  \val{S}=f(a_1)\nimadd f(a_2)\nimadd\cdots\nimadd f(a_n).
\]
Nếu \(\val{S}\ne0\) thì Giáp có chiến lược thắng chắc; nếu không, Ất có
chiến lược thắng chắc. Chứng minh tương tự trò chơi A. Sở dĩ lời giải của
trò chơi B giống trò chơi A là vì quy tắc của hai trò chơi khá gần nhau.

\section{Mở rộng}

\paragraph{Trò chơi C.}
Giáp và Ất đối mặt với một số hàng sỏi. Số sỏi trong mỗi hàng có thể tùy ý.
Ví dụ, vị trí ban đầu có \(n=3\) hàng, với
\[
  a_1=7,\qquad a_2=3,\qquad a_3=3.
\]
Hai người lần lượt lấy sỏi theo các quy tắc sau:
\begin{itemize}
  \item mỗi nước phải lấy đúng hai viên sỏi trong một hàng;
  \item hai viên này phải kề sát nhau;
  \item ai không thể đi đúng quy tắc thì thua.
\end{itemize}

Sơ đồ một vị trí ban đầu:
\[
\begin{array}{c|ccccccc}
&1&2&3&4&5&6&7\\ \hline
\text{hàng thứ nhất }(a_1=7)&\bullet&\bullet&\bullet&\bullet&\bullet&\bullet&\bullet\\
\text{hàng thứ hai }(a_2=3)&\bullet&\bullet&\bullet&&&&\\
\text{hàng thứ ba }(a_3=3)&\bullet&\bullet&\bullet&&&&
\end{array}
\]
Nếu Giáp ở nước đầu chọn lấy hai viên số \(3,4\) trong hàng thứ nhất, thì về
sau không ai có thể lấy cùng lúc hai viên số \(2,5\). Nói cách khác, việc
lấy hai viên \(3,4\) tương đương với tách hàng thứ nhất thành hai hàng có
lần lượt \(2\) và \(3\) viên.

Ta vẫn chỉ quan tâm xem với một vị trí ban đầu, người đi trước hay người đi
sau có chiến lược thắng chắc. Quy tắc của trò chơi C không giống trò chơi A
nhiều như trò chơi B. Tuy vậy, các tính chất then chốt đã liệt kê cho trò
chơi A đều vẫn xuất hiện ở trò chơi C. Ví dụ, vị trí trên được biểu diễn bởi
\((7,3,3)\), và thắng thua của nó giống vị trí \((7)\). Lời giải trò chơi A
được rút ra từ các tính chất của nó. Vì vậy ta phỏng đoán trò chơi C cũng có
thể giải bằng một phương pháp tương tự.

\section{Phần cốt lõi}

Nhắc lại kết luận của trò chơi A và phần mở rộng sang trò chơi B. Với trò
chơi C, ý tưởng là thiết kế một hàm \(f\), xem giá trị của \(f\) như số nhị
phân. Với một vị trí bất kỳ
\[
  S=(a_1,a_2,\ldots,a_n),
\]
đặt
\[
  \val{S}=f(a_1)\nimadd f(a_2)\nimadd\cdots\nimadd f(a_n).
\]
Nếu \(\val{S}\ne0\) thì người đi trước, Giáp, có chiến lược thắng chắc; nếu
\(\val{S}=0\) thì người đi sau, Ất, có chiến lược thắng chắc. Trong trò chơi
A, \(f(x)=x\). Trong trò chơi B, \(f(x)=x\bmod(m+1)\). Vậy trong trò chơi C,
\(f(x)\) là gì?

Điểm then chốt là xây dựng một hàm \(f\) thỏa yêu cầu. Nhìn lại chứng minh
cho trò chơi A và B, việc \(f\) có thỏa yêu cầu hay không phụ thuộc vào việc
\(\val{S}\) có hai tính chất sau:
\begin{itemize}
  \item nếu \(\val{S}=0\), thì với mọi nước đi \(S\to T\), ta có
  \(\val{T}\ne0\);
  \item nếu \(\val{S}\ne0\), thì tồn tại một nước đi \(S\to T\) sao cho
  \(\val{T}=0\).
\end{itemize}

Dùng ký hiệu \(\moves{x}\) để chỉ tập mọi vị trí có thể xuất hiện sau một
nước đi từ vị trí một hàng \((x)\). Ví dụ:
\[
\begin{aligned}
  \moves{3}&=\{(2),(1),(0)\} &&\text{trong trò chơi A},\\
  \moves{9}&=\{(8),(7),(6),(5)\},\quad
  \moves{2}=\{(1),(0)\} &&\text{trong trò chơi B nếu }m=4,\\
  \moves{7}&=\{(5),(1,4),(2,3)\} &&\text{trong trò chơi C}.
\end{aligned}
\]
Định nghĩa tập \(g(x)\): nếu
\[
  \moves{x}=\{S_1,S_2,\ldots,S_k\},
\]
thì
\[
  g(x)=\{\val{S_1},\val{S_2},\ldots,\val{S_k}\}.
\]
Chẳng hạn, trong trò chơi A,
\[
  \moves{3}=\{(2),(1),(0)\}
\]
nên
\[
  g(3)=\{\val{(2)},\val{(1)},\val{(0)}\}=\{10,01,00\}.
\]
Trong trò chơi C,
\[
  g(7)=\{\val{(5)},\val{(1,4)},\val{(2,3)}\}.
\]

Xét điều kiện thứ nhất. Giả sử
\[
  S=(a_1,a_2,\ldots,a_n)
\]
và người đi trước chỉ chọn hàng \(a_1\). Vì
\[
  \val{S}=f(a_1)\nimadd\val{(a_2,\ldots,a_n)}=0,
\]
ta có
\[
  f(a_1)=\val{(a_2,\ldots,a_n)}.
\]
Người đi trước có thể biến vị trí một hàng \((a_1)\) thành một vị trí
\((b_1,\ldots,b_m)\), trong đó
\[
  \val{(b_1,\ldots,b_m)}\in g(a_1).
\]
Khi đó
\[
  T=(b_1,\ldots,b_m)+(a_2,\ldots,a_n),
\]
và
\[
\begin{aligned}
  \val{T}
  &=\val{(b_1,\ldots,b_m)}\nimadd\val{(a_2,\ldots,a_n)}\\
  &=\val{(b_1,\ldots,b_m)}\nimadd f(a_1).
\end{aligned}
\]
Muốn luôn có \(\val{T}\ne0\), điều cần và đủ là
\[
  \val{(b_1,\ldots,b_m)}\ne f(a_1)
\]
với mọi nước đi, tức \(f(a_1)\notin g(a_1)\).

Xét điều kiện thứ hai. Giả sử
\[
  p=\val{S}=f(a_1)\nimadd f(a_2)\nimadd\cdots\nimadd f(a_n)\ne0.
\]
Theo tính chất xor đã chứng minh, tồn tại \(k\) sao cho
\[
  f(a_k)\nimadd p<f(a_k).
\]
Không mất tính tổng quát, lấy \(k=1\) và đặt
\[
  x=f(a_1)\nimadd p.
\]
Vì
\[
  p=f(a_1)\nimadd\val{(a_2,\ldots,a_n)},
\]
ta có
\[
  \val{(a_2,\ldots,a_n)}=p\nimadd f(a_1)=x.
\]
Nếu người đi trước biến \((a_1)\) thành \((b_1,\ldots,b_m)\), vị trí mới là
\[
  T=(b_1,\ldots,b_m)+(a_2,\ldots,a_n),
\]
và
\[
  \val{T}=\val{(b_1,\ldots,b_m)}\nimadd x.
\]
Để \(\val{T}=0\), ta cần tìm \((b_1,\ldots,b_m)\) sao cho
\[
  \val{(b_1,\ldots,b_m)}=x.
\]
Nếu bảo đảm được \(x\in g(a_1)\), ta tìm được nước đi tương ứng. Vì
\(x<f(a_1)\), chỉ cần \(g(a_1)\) chứa toàn bộ tập
\[
  \{0,1,\ldots,f(a_1)-1\}.
\]

Vậy một điều kiện đủ để hàm \(f\) thỏa yêu cầu là:
\begin{itemize}
  \item \(f(a_1)\notin g(a_1)\);
  \item \(g(a_1)\) chứa \(\{0,1,\ldots,f(a_1)-1\}\).
\end{itemize}
Ví dụ, nếu
\[
  g(a_1)=\{0,1,2,5,7,8,9\},
\]
thì \(f(a_1)=3\) thỏa yêu cầu.

Dùng \(\Nat=\{0,1,2,\ldots\}\) cho tập các số nguyên không âm. Lấy \(\Nat\)
làm tập nền và gọi \(G(x)\) là phần bù của \(g(x)\). Định nghĩa
\[
  f(n)=\min G(n).
\]
Nói cách khác, \(f(n)\) là số nhỏ nhất trong \(G(n)\), tức số nguyên không âm
nhỏ nhất không xuất hiện trong \(g(n)\). Với
\[
  S=(a_1,a_2,\ldots,a_n),
\]
đặt
\[
  \val{S}=f(a_1)\nimadd f(a_2)\nimadd\cdots\nimadd f(a_n).
\]
Khi đó
\[
  \val{S}=0 \Longleftrightarrow S\text{ thua},\qquad
  \val{S}\ne0 \Longleftrightarrow S\text{ thắng}.
\]

Các giá trị \(f\) đầu tiên của trò chơi C là:
\begin{itemize}
  \item \(g(0)=\varnothing\), \(G(0)=\{0,1,\ldots\}\), nên \(f(0)=0\);
  \item \(g(1)=\varnothing\), \(G(1)=\{0,1,\ldots\}\), nên \(f(1)=0\);
  \item \(g(2)=\{\val{(0)}\}=\{f(0)\}=\{0\}\), nên \(G(2)=\{1,2,\ldots\}\)
  và \(f(2)=1\);
  \item \(g(3)=\{\val{(1)}\}=\{f(1)\}=\{0\}\), nên \(G(3)=\{1,2,\ldots\}\)
  và \(f(3)=1\);
  \item \(g(4)=\{\val{(2)},\val{(1,1)}\}=\{f(2),f(1)\nimadd f(1)\}=\{1,0\}\),
  nên \(G(4)=\{2,3,\ldots\}\) và \(f(4)=2\);
  \item \(g(5)=\{\val{(3)},\val{(1,2)}\}=\{f(3),f(1)\nimadd f(2)\}=\{1\}\),
  nên \(G(5)=\{0,2,3,\ldots\}\) và \(f(5)=0\);
  \item dòng PDF ghi
  \(g(6)=\{\val{(4)},\val{(1,4)},\val{(2,2)}\}=\{2,1,0\}\), nên
  \(G(6)=\{3,4,\ldots\}\) và \(f(6)=3\). Giá trị \(1\) ở giữa phù hợp với
  \(\val{(1,3)}\), vì vậy \((1,4)\) nhiều khả năng là lỗi đánh máy hoặc lỗi
  nhận dạng trong PDF.
  \item \(g(7)=\{\val{(5)},\val{(1,4)},\val{(2,3)}\}=\{0,2,0\}\), nên
  \(G(7)=\{1,3,4,\ldots\}\) và \(f(7)=1\).
\end{itemize}

Với vị trí trong ví dụ
\[
  S=(7,3,3),
\]
ta có
\[
  \val{S}=f(7)\nimadd f(3)\nimadd f(3)=1\nimadd1\nimadd1=1,
\]
nên \(S\) thắng. Với vị trí
\[
  S=(3,4,6),
\]
theo bảng giá trị trong văn bản gốc,
\[
  \val{S}=1\nimadd2\nimadd3=01\nimadd10\nimadd11=0,
\]
nên \(S\) thua.

\section{Kết luận}

Lời giải tổng quát cho lớp trò chơi này là:
\begin{enumerate}
  \item dùng một bộ \(n\) phần tử \((a_1,a_2,\ldots,a_n)\) để mô tả một vị
  trí;
  \item dùng \(\val{S}\) để biểu diễn số nhị phân ứng với vị trí \(S\);
  \item dùng \(\moves{x}\) để biểu diễn tập mọi vị trí có thể xuất hiện sau
  một nước đi từ vị trí một hàng \((x)\);
  \item nếu \(\moves{x}=\{S_1,S_2,\ldots,S_k\}\), định nghĩa
  \(g(x)=\{\val{S_1},\val{S_2},\ldots,\val{S_k}\}\);
  \item lấy \(\Nat\) làm tập nền và gọi \(G(x)\) là phần bù của \(g(x)\);
  \item định nghĩa \(f(n)=\min G(n)\);
  \item với \(S=(a_1,a_2,\ldots,a_n)\), đặt
  \[
    \val{S}=f(a_1)\nimadd f(a_2)\nimadd\cdots\nimadd f(a_n),
  \]
  trong đó \(\nimadd\) là phép cộng nhị phân không nhớ.
\end{enumerate}
Nếu \(\val{S}\ne0\) thì người đi trước có chiến lược thắng chắc; nếu
\(\val{S}=0\) thì người đi sau có chiến lược thắng chắc.

Phạm vi áp dụng và các điều kiện ràng buộc:
\begin{itemize}
  \item trò chơi lấy sỏi hai người và các trò chơi tương tự;
  \item mỗi nước chỉ được thao tác trên một đống hoặc một hàng sỏi;
  \item ràng buộc của mỗi nước chỉ phụ thuộc vào số sỏi trong đống đó hoặc
  vào một vài hằng số;
  \item quá trình chơi kết thúc sau hữu hạn bước, không sinh chu trình;
  \item ai không thể tiếp tục đi thì thua.
\end{itemize}

\paragraph{Trò chơi D, POI 2000, Stripes.}
Một hàng có \(L\) viên sỏi. Giáp và Ất lần lượt lấy \(A\), \(B\), hoặc \(C\)
viên sỏi kề sát nhau. Ai không thể lấy thì thua. Biết \(A,B,C,L\), hỏi ai có
chiến lược thắng chắc. Với kết luận phía trên, trò chơi này không còn khó.

\section{Tổng kết}

\subsection{Từ góc nhìn tối ưu thuật toán}

Trò chơi lấy sỏi thuộc một lớp trò chơi đối kháng điển hình. Về lý thuyết,
liệt kê mọi vị trí có thể cho ta chiến lược tối ưu. Nhưng độ phức tạp thời
gian và không gian của việc liệt kê quá cao. Lời giải trong bài kiểm soát
hiệu quả độ phức tạp và có thể xem là một cách tối ưu phương pháp liệt kê.

Quá trình tối ưu thuật toán có thể xem là quá trình tối ưu cách biểu diễn vị
trí. Ban đầu, ta dùng bộ \(n\) phần tử để biểu diễn một vị trí; đây là cách
rất trực quan. Vì ta chỉ quan tâm thắng thua, ta thu được tính chất đầu tiên:
bộ \(n\) phần tử không có thứ tự. Phân tích sâu hơn cho thấy nếu trong bộ có
hai số giống nhau, ta bỏ cặp đó đi mà không ảnh hưởng đến thắng thua. Vì vậy
ta chuyển sang dùng tập hợp để biểu diễn vị trí. Cuối cùng, qua so sánh với
số nhị phân, ta đơn giản hóa tới mức dùng một số để biểu diễn một vị trí.

Tối ưu cách biểu diễn vị trí làm giảm rất nhiều lượng tìm kiếm. Phần giảm đó
đã đi đâu? Xét các vị trí trong trò chơi A:
\[
  (3,3,1),\quad (1,3,3),\quad (5,5,1),\quad (2,3).
\]
\begin{enumerate}
  \item Dùng bộ \(n\) phần tử: các vị trí này được xem là khác nhau.
  \item Dùng bộ \(n\) phần tử không thứ tự: \((3,3,1)\) và \((1,3,3)\) là
  cùng một vị trí.
  \item Dùng tập hợp: \((3,3,1)\), \((1,3,3)\), và \((5,5,1)\) đều giống
  nhau, có thể biểu diễn bằng \(\{1\}\).
  \item Dùng số nhị phân: cả bốn vị trí trên đều có giá trị nhị phân bằng
  \(1\), nên được xem là cùng loại theo thắng thua.
\end{enumerate}

Bản chất của tối ưu thuật toán là tránh liệt kê lặp lại các vị trí giống
nhau. Điểm then chốt nằm ở định nghĩa ``vị trí giống nhau''. Định nghĩa ấy
phải phản ánh được tính chất của trò chơi. Đây là lý do ta không đơn giản
phân loại vị trí chỉ bằng kết quả thắng thua.

\subsection{Từ góc nhìn xây dựng thuật toán}

Trong quá trình nhận thức sự vật, ban đầu con người chỉ thấy các hiện tượng
riêng lẻ. Đó là giai đoạn nhận thức cảm tính. Ở giai đoạn này, ta chưa thể
đưa ra kết luận hợp logic. Khi nghiên cứu đi sâu hơn, những cảm giác và ấn
tượng ấy lặp lại nhiều lần, tạo ra một bước chuyển trong nhận thức, rồi cuối
cùng sinh ra kết luận hợp logic. Đó là giai đoạn nhận thức lý tính.

Quá trình nhận thức sự vật là quá trình đi từ nhận thức cảm tính lên nhận
thức lý tính. Cụ thể với việc giải lớp trò chơi này, ta phải bắt đầu từ điều
đơn giản. Khi gặp một vấn đề phức tạp, có lẽ ai cũng biết nên bắt đầu từ
trường hợp đơn giản, nhưng không phải ai cũng rút ra được quy luật tổng quát
từ đó. Vậy ta đi từ nông đến sâu bằng cách nào?

Hai đống sỏi có cùng số lượng là một vị trí rất đơn giản. Ta bắt đầu ở đó,
so sánh một đống sỏi với một vị trí con, rồi rút ra kết luận khi hai vị trí
con bằng nhau. Trên cơ sở đó, ta nghiên cứu quan hệ giữa thắng thua của một
vị trí với các vị trí con của nó, và kết luận rằng có thể dùng tập hợp để mô
tả một vị trí. Nhưng ta không dừng lại ở bước này. Ta tiếp tục so sánh phép
phân tích vị trí với phép cộng nhị phân, từ đó phát hiện quan hệ giữa vị trí
và số nhị phân. Quá trình này gọi là \term{từ cái này suy sang cái kia}.

Khi phân tích kết luận ``dùng tập hợp để biểu diễn vị trí'', ta thấy về bản
chất nó đã đơn giản hóa cách biểu diễn vị trí. Từ đó ta liên tưởng liệu có
thể đơn giản hóa hơn nữa, chẳng hạn dùng một số để biểu diễn vị trí. Khi
giải trò chơi C, ta không đặt nặng quy tắc của nó khác trò chơi A nhiều hay
ít, mà chú ý rằng nó có các tính chất tương tự trò chơi A, rồi nghĩ tới dùng
phương pháp tương tự để giải. Quá trình này gọi là \term{từ biểu hiện đi vào
bản chất}.

Trong quá trình giải trò chơi A và B, ta tích lũy được nhiều kinh nghiệm.
Nhưng khi giải trò chơi C, ta chỉ nhắc tới phần tinh túy của trò chơi A và B:
xây dựng một hàm \(f\). Đó là \term{gạn phần thô, giữ phần tinh}.

Khi so sánh vị trí với số nhị phân, ban đầu ta thử gắn trực tiếp thắng thua
của vị trí với \(1\) và \(0\) trong nhị phân. Sau khi thấy không ổn, ta đổi
sang so sánh với toàn bộ số nhị phân. Bước này gọi là \term{bỏ điều sai, giữ
điều đúng}.

``Từ cái này suy sang cái kia, từ biểu hiện đi vào bản chất, gạn phần thô,
giữ phần tinh, bỏ điều sai, giữ điều đúng'' chính là điểm then chốt để đi từ
nhận thức cảm tính lên nhận thức lý tính.

\appendix

\section{Phụ lục: nội dung trình chiếu trích xuất từ PDF}

Phần sau của PDF chứa một bộ trang trình chiếu đi kèm bài viết. Nội dung
dưới đây giữ lại các ý và ví dụ trích xuất được, chuyển sang văn bản LaTeX.

\subsection{Quá trình nhận thức sự vật}

Từ hiện tượng của sự vật đi tới bản chất của sự vật là quá trình chuyển từ
nhận thức cảm tính sang nhận thức lý tính. Con người nhận thức sự vật luôn
bắt đầu từ điều đơn giản. Nhưng không phải ai cũng có thể nhận ra quy luật
tổng quát từ các sự vật đơn giản. Vấn đề là làm thế nào để đi từ nông đến
sâu.

\subsection{Trò chơi trong phần trình chiếu}

Hai người đối mặt với một số hàng sỏi, số sỏi mỗi hàng tùy ý. Mỗi nước phải
lấy đúng hai viên kề sát nhau trong một hàng; ai không thể lấy thì thua. Ví
dụ:
\[
\begin{array}{c|ccccccc}
&1&2&3&4&5&6&7\\ \hline
a_1=7&\bullet&\bullet&\bullet&\bullet&\bullet&\bullet&\bullet\\
a_2=3&\bullet&\bullet&\bullet&&&&\\
a_3=5&\bullet&\bullet&\bullet&\bullet&\bullet&&
\end{array}
\]
Nếu một hàng có \(7\) viên và ta lấy viên \(3,4\), hàng ấy được xem như tách
thành hai hàng có \(2\) và \(3\) viên. Vì vậy số hàng của vị trí có thể tăng
lên trong quá trình chơi.

\subsection{Bắt đầu từ trường hợp đơn giản trong phần trình chiếu}

Dùng một bộ không thứ tự \((a_1,a_2,\ldots,a_n)\) để mô tả một vị trí; ví dụ
\((3,3,1)\). Nếu người đi trước có chiến lược thắng chắc, gọi đó là vị trí
thắng; nếu người đi sau có chiến lược thắng chắc, gọi đó là vị trí thua. Nếu
vị trí ban đầu phân thành hai vị trí con giống nhau thì Ất có chiến lược
thắng chắc:
\[
  X+X=\text{thua}.
\]
Các quan hệ được nêu trên trang trình chiếu là
\[
\begin{array}{rcl}
X+X&=&\text{thua},\\
\text{thắng}+\text{thua}&=&\text{thắng},\\
\text{thua}+\text{thua}&=&\text{thua},\\
\text{thắng}+\text{thắng}&=&\text{chưa xác định},\\
X+\text{thua}&=&X,\\
X+C+C&=&X.
\end{array}
\]

\subsection{Vị trí và tập hợp}

Vì chỉ quan tâm thắng thua, một vị trí có thể được mô tả bằng một tập hợp.
Ví dụ,
\[
  (2,2,2,7,9,9)=(2,2,2,7)+(9)+(9)
\]
có cùng thắng thua với \((2,2,2,7)\), rồi với \((2,7)\), nên có thể biểu
diễn bằng \(\{2,7\}\). Bản chất ở đây là đơn giản hóa cách biểu diễn vị trí.
Câu hỏi tiếp theo là có thể đơn giản hóa hơn nữa không.

\subsection{Tương tự trong phần trình chiếu}

Các trang trình chiếu so sánh phép cộng vị trí với phép cộng \(0,1\) trong
nhị phân:
\[
\begin{array}{rcl@{\qquad}rcl}
\text{thắng}+\text{thua}&=&\text{thắng},&
1\nimadd0&=&1,\\
\text{thua}+\text{thắng}&=&\text{thắng},&
0\nimadd1&=&1,\\
\text{thua}+\text{thua}&=&\text{thua},&
0\nimadd0&=&0,\\
\text{thắng}+\text{thắng}&=&\text{chưa xác định},&
1\nimadd1&=&0.
\end{array}
\]
Vì \(1\nimadd1=0\) không khớp hoàn toàn với trường hợp thắng cộng thắng, ta
chuyển sang so sánh với số nhị phân nói chung: thắng tương ứng với giá trị
khác \(0\), thua tương ứng với giá trị \(0\).

\subsection{Liên tưởng trong phần trình chiếu}

Dùng \(\val{S}\) để chỉ giá trị của vị trí \(S\). Khi đó
\[
  \val{S}=0 \Longleftrightarrow S\text{ thua},\qquad
  \val{S}\ne0 \Longleftrightarrow S\text{ thắng},
\]
và
\[
  \val{(a_1,\ldots,a_n)}
  =f(a_1)\nimadd\cdots\nimadd f(a_n).
\]
Với ví dụ:
\[
\begin{aligned}
  (3,3)&=(3)+(3),&
  \val{(3,3)}&=\val{(3)}\nimadd\val{(3)},\\
  (3,3,1)&=(3)+(3)+(1),&
  \val{(3,3,1)}&=f(3)\nimadd f(3)\nimadd f(1).
\end{aligned}
\]
Điểm then chốt là xây dựng hàm \(f(x)\).

\subsection{Cách xây dựng trong phần trình chiếu}

Tập \(g(x)\) biểu diễn các giá trị của mọi vị trí có thể nhận được sau một
nước đi từ \((x)\). Ví dụ, khi lấy hai viên kề nhau:
\[
  (7)\to (5),\quad (7)\to(1,4),\quad (7)\to(2,3),
\]
nên
\[
  g(7)=\{\val{(5)},\val{(1,4)},\val{(2,3)}\}.
\]
Có thể chứng minh rằng chọn \(f(x)\) là số nguyên không âm nhỏ nhất chưa xuất
hiện trong \(g(x)\) thì thỏa yêu cầu. Nếu \(G(x)\) là phần bù của \(g(x)\)
trong tập các số nguyên không âm, thì
\[
  f(x)=\min G(x).
\]
Ví dụ, nếu
\[
  g(x)=\{0,1,2,5,7,8,9\},
\]
thì \(f(x)=3\).

Một bảng ví dụ trên trang trình chiếu:
\[
\begin{array}{c|cccccccc}
x&0&1&2&3&4&5&6&7\\ \hline
f(x)&0&0&1&1&2&0&3&?
\end{array}
\]
Với \(x=7\), các vị trí kế tiếp được liệt kê:
\[
\begin{array}{rcl}
(5)&:&\val{(5)}=f(5)=0,\\
(1,4)&:&\val{(1,4)}=f(1)\nimadd f(4)=0\nimadd2=2,\\
(2,3)&:&\val{(2,3)}=f(2)\nimadd f(3)=1\nimadd1=0.
\end{array}
\]
Do đó
\[
  g(7)=\{0,2\},\qquad G(7)=\{1,3,4,5,\ldots\},
\]
và
\[
  f(7)=\min G(7)=1.
\]
Ví dụ vị trí
\[
  (7,3,5)
\]
có
\[
  \val{(7,3,5)}=f(7)\nimadd f(3)\nimadd f(5)=1\nimadd1\nimadd0=0,
\]
nên là vị trí thua; người đi sau, Ất, có chiến lược thắng chắc.

\subsection{Mở rộng trong phần trình chiếu}

Ta có thể thay đổi quy tắc:
\begin{itemize}
  \item mỗi lần lấy hai viên kề sát nhau;
  \item mỗi lần lấy ba viên kề sát nhau;
  \item mỗi lần lấy một số viên kề sát nhau tùy ý;
  \item mỗi lần lấy hai viên bất kỳ trong một hàng, không yêu cầu kề nhau;
  \item mỗi lần lấy một số viên bất kỳ trong một hàng, không yêu cầu kề nhau.
\end{itemize}
Đặc điểm chung của lớp trò chơi này:
\begin{itemize}
  \item hai người lấy sỏi;
  \item mỗi nước chỉ thao tác trên một hàng sỏi;
  \item ràng buộc của mỗi nước chỉ phụ thuộc vào số sỏi trong hàng đó hoặc
  vào một số hằng số;
  \item quá trình chơi kết thúc sau hữu hạn bước, không sinh chu trình;
  \item ai không thể tiếp tục đi thì thua.
\end{itemize}

\subsection{Lời giải tổng quát trong phần trình chiếu}

Để phán đoán ai có chiến lược thắng chắc:
\begin{itemize}
  \item đặt vị trí \(S=(a_1,a_2,\ldots,a_n)\);
  \item giá trị của \(S\) là
  \[
    \val{S}=f(a_1)\nimadd\cdots\nimadd f(a_n);
  \]
  \item nếu \(\val{S}\ne0\), người đi trước thắng chắc;
  \item nếu \(\val{S}=0\), người đi sau thắng chắc.
\end{itemize}
Cách tìm hàm \(f(x)\):
\begin{itemize}
  \item \(f(0)=0\);
  \item \(g(x)\) là tập giá trị của các vị trí có thể nhận được sau một nước
  đi từ \((x)\);
  \item \(G(x)\) là phần bù của \(g(x)\) trong tập các số nguyên không âm;
  \item \(f(x)=\min G(x)\).
\end{itemize}

\subsection{Các trang tổng kết trong phần trình chiếu}

\paragraph{Ưu điểm và hạn chế.}
Ưu điểm: phạm vi áp dụng rộng, có thể dùng trực tiếp cho đa số trò chơi thuộc
lớp này; so với liệt kê, tốc độ nhanh hơn và độ phức tạp thời gian, không
gian thấp hơn. Hạn chế: với một số trò chơi, có thể suy ra trực tiếp công
thức của hàm \(f(x)\). Ví dụ, trò chơi có nhiều đống sỏi, mỗi lượt lấy trong
một đống ít nhất một viên và không giới hạn tối đa, ai không lấy được thì
thua; với trò chơi này, có thể chứng minh \(f(x)=x\) đã thỏa yêu cầu.

\paragraph{Tối ưu thuật toán.}
Cách làm trên có thể xem là tối ưu thuật toán tìm kiếm. Quá trình tối ưu
thuật toán có thể xem là quá trình đơn giản hóa cách biểu diễn vị trí. Bản
chất là tránh liệt kê lặp lại các vị trí giống nhau. Ví dụ:
\[
\begin{array}{c|ccccc}
\text{bộ không thứ tự}
&(2,5,5)&(5,2,5)&(2,3,3)&(2,3,4,6)&(3,4)\\ \hline
\text{tập hợp}
&\{2\}&\{2\}&\{2\}&\{2,3,4,6\}&\{3,4\}\\ \hline
\text{số nhị phân}
&01&01&01&01&11
\end{array}
\]

\paragraph{Đi từ nông đến sâu.}
Các trang cuối gom lại các bước:
\[
\begin{array}{c}
X+X=\text{thua},\quad
\text{thắng}+\text{thua}=\text{thắng},\quad
\text{thua}+\text{thua}=\text{thua},\\
\text{thắng}+\text{thắng}=\text{chưa xác định},\quad
X+C+C=X.
\end{array}
\]
Từ đó đi qua các ý: dùng \(0\) và \(1\) để biểu diễn thắng thua, rồi sửa
thành dùng số nhị phân để biểu diễn vị trí; dùng tập hợp để biểu diễn vị
trí, rồi nhận ra bản chất là đơn giản hóa cách biểu diễn; cuối cùng đạt tới
\[
  f(x)=\min G(x).
\]
Con đường từ nhận thức cảm tính đến nhận thức lý tính gồm: bỏ điều sai, giữ
điều đúng; gạn phần thô, giữ phần tinh; từ cái này suy sang cái kia; từ biểu
hiện đi vào bản chất.

\paragraph{Kết thúc.}
Ba kết luận chính của bộ trang trình chiếu là:
\begin{itemize}
  \item lời giải tổng quát của lớp trò chơi này là \(f(x)=\min G(x)\);
  \item bản chất của tối ưu thuật toán là tránh tìm kiếm lặp lại;
  \item cách đi từ nông đến sâu là bỏ điều sai, giữ điều đúng; gạn phần thô,
  giữ phần tinh; từ cái này suy sang cái kia; từ biểu hiện đi vào bản chất.
\end{itemize}

\end{document}
